% !TeX program = xelatex
\documentclass[a4paper,10pt]{article}

% ============================================================================
% សន្លឹកកិច្ចការសិស្ស ស-ខ្មៅ (Student Worksheet - Grade 8 3.1.1 PROP1)
% បរិបទ៖ គ្មានស្លាយ LCD និងគ្មានម៉ាស៊ីនគិតលេខ — សិស្សបំពេញផ្ទាល់លើសន្លឹកកិច្ចការ ២ ទំព័រ A4
% ============================================================================
\usepackage[a4paper, left=0.85cm, right=0.85cm, top=0.70cm, bottom=0.95cm, footskip=0.48cm]{geometry}
\usepackage{amsmath, amssymb}
\usepackage{unicode-math}
\usepackage{fontspec}
\usepackage{array, colortbl}
\usepackage[table]{xcolor}
\usepackage{enumitem}
\usepackage{pifont}
\usepackage{fancyhdr}
\usepackage{lastpage}
\usepackage{tikz}
\usepackage[most]{tcolorbox}
\usetikzlibrary{calc, arrows.meta, decorations.pathreplacing, positioning, babel}

\makeatletter
\let\@expl@@@mark@update@singlecol@structures@@\relax
\makeatother

\providecommand{\NGSFindFonts}{%
  \IfFileExists{C:/Windows/Fonts/cambria.ttc}{%
    \def\NGSSysPath{C:/Windows/Fonts/}%
    \def\NGSUserPath{C:/Users/KIT/AppData/Local/Microsoft/Windows/Fonts/}%
  }{%
    \IfFileExists{./fonts/KhmerOSsiemreap.ttf}{\def\NGSSysPath{./fonts/}\def\NGSUserPath{./fonts/}}{%
    \IfFileExists{../fonts/KhmerOSsiemreap.ttf}{\def\NGSSysPath{../fonts/}\def\NGSUserPath{../fonts/}}{%
    \IfFileExists{../../fonts/KhmerOSsiemreap.ttf}{\def\NGSSysPath{../../fonts/}\def\NGSUserPath{../../fonts/}}{%
    \IfFileExists{../../../fonts/KhmerOSsiemreap.ttf}{\def\NGSSysPath{../../../fonts/}\def\NGSUserPath{../../../fonts/}}{%
    \IfFileExists{../../../../fonts/KhmerOSsiemreap.ttf}{\def\NGSSysPath{../../../../fonts/}\def\NGSUserPath{../../../../fonts/}}{%
    \IfFileExists{../../../../../fonts/KhmerOSsiemreap.ttf}{\def\NGSSysPath{../../../../../fonts/}\def\NGSUserPath{../../../../../fonts/}}{%
    \IfFileExists{../../../../../../fonts/KhmerOSsiemreap.ttf}{\def\NGSSysPath{../../../../../../fonts/}\def\NGSUserPath{../../../../../../fonts/}}{%
      \def\NGSSysPath{../../../../../../../fonts/}\def\NGSUserPath{../../../../../../../fonts/}%
    }}}}}}}%
  }%
}
\NGSFindFonts

\setmathfont[Path={\NGSSysPath}, FontIndex=1, Scale=1.02]{cambria.ttc}
\defaultfontfeatures{Ligatures=TeX}

\setmainfont[
  Path={\NGSUserPath}, Extension=.ttf, Script=Khmer,
  BoldFont={KhmerOSsiemreap}, BoldFeatures={FakeBold=1.8},
  ItalicFont={KhmerOSsiemreap}, ItalicFeatures={FakeSlant=0.22}
]{KhmerOSsiemreap}

\newfontfamily\muol[Path={\NGSUserPath}, Extension=.ttf, Script=Khmer]{KhmerOSmuollight}
\newfontfamily\TimesEn[Path={\NGSSysPath}, Extension=.ttf, BoldFont={timesbd}, ItalicFont={timesi}]{times}

\XeTeXlinebreaklocale "km"
\XeTeXlinebreakskip = 0pt plus 1pt minus 0.1pt

\definecolor{panelgray}{gray}{0.93}
\definecolor{headgray}{gray}{0.85}

\pagestyle{fancy}
\fancyhf{}
\renewcommand{\headrulewidth}{0pt}
\fancyfoot[L]{\fontsize{7.5pt}{9.5pt}\selectfont អនុវិទ្យាល័យព្រះក្រសាំង \en{|} គណិតវិទ្យាថ្នាក់ទី៨ \en{|} មេរៀនទី៣៖ ១.១. ទំហំសមាមាត្រស្រប (គ្រូ៖ ហុក ប៊ុនថែន)}
\fancyfoot[R]{\fontsize{7.8pt}{9.5pt}\selectfont ទំព័រ \en{\thepage/\pageref{LastPage}}}

\setlength{\parindent}{0pt}
\setlength{\parskip}{0pt}
\setlength{\tabcolsep}{3.8pt}

\newcommand{\en}[1]{{\TimesEn #1}}
\newcommand{\enbf}[1]{{\TimesEn\bfseries #1}}
\newcommand{\unit}[1]{\,\text{\en{#1}}}
\newcommand{\cnum}[1]{\mbox{\fontsize{10pt}{12pt}\selectfont\ding{\numexpr181+#1\relax}}}
\newcommand{\chkbox}{\tikz[baseline=-0.12ex]\draw[black, line width=0.8pt] (0,0) rectangle (0.24cm,0.24cm);}
\newcommand{\blank}[1][1.5cm]{\makebox[#1]{\dotfill}}
\newcommand{\ansbox}[1][0.85cm]{\,\tikz[baseline=-0.5ex]\draw[black, line width=0.65pt, fill=white] (0,0) rectangle (#1, 0.42cm);\,}
\newcommand{\secbar}[2]{%
  \begin{tcolorbox}[enhanced, sharp corners, colback=headgray, colframe=black, boxrule=0.9pt, boxsep=1pt, top=2pt, bottom=2pt, left=5pt, right=5pt, before skip=4pt, after skip=3pt]
    {\muol\fontsize{9.2pt}{13pt}\selectfont #1}\hfill{\fontsize{8pt}{11pt}\selectfont\textbf{#2}}
  \end{tcolorbox}%
}

\newcolumntype{L}[1]{>{\raggedright\arraybackslash}p{#1}}
\newcolumntype{C}[1]{>{\centering\arraybackslash}p{#1}}

\begin{document}
\fontsize{9.6pt}{14.4pt}\selectfont

% ============================================================================
% ទំព័រទី ១៖ E1 + E2 (POE 1 & POE 2) + E3 សំយោគនិយមន័យ និងធនាគារហេតុផល
% ============================================================================
\begin{tcolorbox}[enhanced, sharp corners, colback=white, colframe=black, boxrule=1.0pt, boxsep=2pt, top=3pt, bottom=3pt, left=6pt, right=6pt]
  \begin{center}
    {\muol\fontsize{10.4pt}{15pt}\selectfont អនុវិទ្យាល័យព្រះក្រសាំង \enbf{--} សន្លឹកកិច្ចការសិស្ស \enbf{(}មេរៀនទី៣៖ ១\enbf{.}១\enbf{.} ទំហំសមាមាត្រស្រប\enbf{)}}\\[3pt]
    {\fontsize{8.8pt}{12pt}\selectfont
    ឈ្មោះសិស្ស / ដៃគូ៖ \blank[5.2cm] \hspace{0.25cm}
    ថ្នាក់ទី៨៖ \blank[1.1cm] \hspace{0.25cm}
    ក្រុមទី៖ \blank[1.0cm] \hspace{0.25cm}
    កាលបរិច្ឆេទ៖ ២១/១១/២០២៦}
  \end{center}
\end{tcolorbox}

% --- E1 + E2 ---
\secbar{ផ្នែកទី១៖ បេសកកម្មទស្សន៍ទាយ \enbf{(E1)} និងការរុករកជាដៃគូ \enbf{(E2 -- POE 1 \& POE 2)}}{រយៈពេល ១៤ នាទី}

\begin{tcolorbox}[enhanced, sharp corners, colback=panelgray, colframe=black, boxrule=0.75pt, boxsep=1.5pt, top=2.5pt, bottom=2.5pt, left=5pt, right=5pt]
  \textbf{បេសកកម្ម \en{E1} (គូសធីក \chkbox\ មុនពិសោធន៍)៖}
  (១) រ៉ឺស័រព្យួរ $50\unit{g}\to 10\unit{mm}$, $100\unit{g}\to 20\unit{mm}$ — បើព្យួរ $250\unit{g}$ លូត៖ \chkbox\ $40\unit{mm}$ \quad \chkbox\ $50\unit{mm}$\\
  (២) ប្អូន $5$ ឆ្នាំ បង $10$ ឆ្នាំ ; ពេលប្អូន $10$ ឆ្នាំ បង $15$ ឆ្នាំ — តើអាយុបង-ប្អូនសមាមាត្រស្របឬទេ? \chkbox\ មែន \quad \chkbox\ មិនមែន
\end{tcolorbox}

\vspace{3pt}
\textbf{ក. សកម្មភាពទី ១ (\en{POE 1})៖ តារាងពិសោធន៍ប្រវែងលូតនៃរ៉ឺស័រ (សៀវភៅពុម្ព ទំ.២៧)}

{\renewcommand{\arraystretch}{1.48}
\begin{tabular}{|L{5.60cm}|C{2.25cm}|C{2.25cm}|C{2.25cm}|C{2.25cm}|C{2.80cm}|}
  \hline
  \rowcolor{headgray}
  \textbf{បរិមាណក្នុងពិសោធន៍} & \textbf{លើកទី ១} & \textbf{លើកទី ២} & \textbf{លើកទី ៣} & \textbf{លើកទី ៤} & \textbf{សន្និដ្ឋាន} \\
  \hline
  ទម្ងន់ព្យួរនឹងរ៉ឺស័រ $x$ ($\unit{g}$) & $50$ & $100$ & $150$ & $250$ & កើនតាមពហុគុណ \\
  \hline
  ប្រវែងសាច់លូតនៃរ៉ឺស័រ $y$ ($\unit{mm}$) & $10$ & $20$ & $30$ & $\blank[1.2cm]$ & កើនតាមពហុគុណដូចគ្នា \\
  \hline
  ផលធៀប $\dfrac{y}{x}$ (ជាទសភាគ) & $\frac{10}{50} = 0.2$ & $\frac{20}{100} = \blank[0.9cm]$ & $\frac{30}{150} = \blank[0.9cm]$ & $\frac{y}{250} = \blank[0.9cm]$ & ថេរស្មើ $a = \blank[1.0cm]$ \\
  \hline
\end{tabular}}

\vspace{3pt}
$\bullet$ \textbf{ពន្យល់ (\en{E})៖} ដោយផលធៀប $\dfrac{y}{x} = 0.2$ ថេរជានិច្ច យើងសរសេរសមីការទំនាក់ទំនងបាន៖ $y = \blank[1.4cm]\,x$ ។

\vspace{4pt}
\textbf{ខ. សកម្មភាពទី ២ (\en{POE 2})៖ បំបែកអន្ទាក់ «កើនទាំងពីរ» — ប្រៀបធៀបតារាង ក និងតារាង ខ}

\begin{tabular}{|L{9.20cm}|L{9.20cm}|}
  \hline
  \textbf{តារាង ក៖ ប្រាក់ពិន័យបណ្ណាល័យ ($100\text{៛}$/ថ្ងៃ, ទំ.២៧)}\par\vspace{2pt}
  $\bullet$ ចំនួនថ្ងៃ $x$៖ $1,\; 2,\; 5$ $\to$ ប្រាក់ពិន័យ $y$៖ $100,\; 200,\; 500$\par\vspace{2pt}
  $\bullet$ គណនាផលធៀប៖ $\frac{100}{1} = 100$ ; $\frac{200}{2} = \blank[0.9cm]$ ; $\frac{500}{5} = \blank[0.9cm]$\par\vspace{2pt}
  $\bullet$ តើ $\frac{y}{x}$ ថេរឬទេ? \chkbox\ ថេរ \quad \chkbox\ មិនថេរ\par\vspace{2pt}
  $\bullet$ សន្និដ្ឋាន៖ \chkbox\ សមាមាត្រស្រប \quad \chkbox\ មិនសមាមាត្រស្រប
  &
  \textbf{តារាង ខ៖ អាយុប្អូន ($x$) និងអាយុបង ($y = x + 5$)}\par\vspace{2pt}
  $\bullet$ អាយុប្អូន $x$៖ $5,\; 10,\; 15$ $\to$ អាយុបង $y$៖ $10,\; 15,\; 20$\par\vspace{2pt}
  $\bullet$ គណនាផលធៀប៖ $\frac{10}{5} = 2$ ; $\frac{15}{10} = \blank[0.9cm]$ ; $\frac{20}{15} \approx \blank[0.9cm]$\par\vspace{2pt}
  $\bullet$ តើ $\frac{y}{x}$ ថេរឬទេ? \chkbox\ ថេរ \quad \chkbox\ មិនថេរ\par\vspace{2pt}
  $\bullet$ សន្និដ្ឋាន៖ \chkbox\ សមាមាត្រស្រប \quad \chkbox\ មិនសមាមាត្រស្រប \\
  \hline
\end{tabular}

\vspace{4pt}
$\bullet$ \textbf{ពន្យល់ (\en{E})៖} ហេតុអ្វីក្នុងតារាង ខ អាយុប្អូន និងអាយុបងកើន $5$ ឆ្នាំដូចគ្នា តែមិនមែនជាសមាមាត្រស្រប? ព្រោះទំនាក់ទំនងនេះជាការ\textbf{បូកថេរ} ($y = x + \blank[0.8cm]$) ធ្វើឱ្យផលធៀប $\dfrac{y}{x}$ \blank[2.0cm] ។

% --- E3 ---
\secbar{ផ្នែកទី២៖ សំយោគក្បួនទូទៅ \enbf{(E3)} និងធនាគារហេតុផល \enbf{B\&W} \cnum{1}\enbf{--}\cnum{5}}{សៀវភៅពុម្ព ទំ.២៧}

\begin{tabular}{|L{9.20cm}|L{9.20cm}|}
  \hline
  \textbf{១. ក្បួនជាទូទៅ (ទំ.២៧)៖}\par\vspace{2pt}
  $\bullet$ ប្រសិនបើទំហំ $y$ \textbf{សមាមាត្រស្រប}នឹងទំហំ $x$ គេបាន៖\par\vspace{3pt}
  {\centering $\dfrac{y}{x} = \ansbox[0.9cm]$ \quad ឬ \quad $y = \ansbox[0.9cm]\,x$ \quad ($a \neq 0$)\par}\vspace{3pt}
  $\bullet$ ចំនួនថេរ $a$ ហៅថា \textbf{មេគុណសមាមាត្រ} (តម្លៃក្នុង $1$ ឯកតា) ។\par\vspace{2pt}
  $\bullet$ ពេល $x = 0$ នោះ $y = \blank[0.8cm]$ ជានិច្ច ។
  &
  \textbf{២. ធនាគារហេតុផល \en{B\&W} (\cnum{1}--\cnum{5})៖}\par\vspace{2pt}
  \cnum{1} និយមន័យ៖ $y$ សមាមាត្រស្របនឹង $x \iff \frac{y}{x} = a$ ថេរ ($y=ax$) ។\par\vspace{2pt}
  \cnum{2} មេគុណសមាមាត្រ៖ $a = \frac{y}{x}$ ជាតម្លៃនៃ $y$ ក្នុង $1$ ឯកតានៃ $x$ ។\par\vspace{2pt}
  \cnum{3} ពហុគុណ៖ បើ $x$ គុណនឹង $k$ នោះ $y$ ក៏គុណនឹង $k$ ដែរ ។\par\vspace{2pt}
  \cnum{4} អន្ទាក់បូកថេរ៖ $y = x + c$ ($c\neq 0$) មិនមែនសមាមាត្រស្របទេ! \par\vspace{2pt}
  \cnum{5} គល់ $(0,0)$៖ បើ $x=0$ ហើយ $y\neq 0$ នោះមិនមែនសមាមាត្រស្របទេ ។ \\
  \hline
\end{tabular}

\vspace{5pt}
\textbf{៣. លំហាត់ត្រួតពិនិត្យរហ័ស៖ កំណត់មេគុណសមាមាត្រ $a$ និងសមីការ $y = ax$}

{\renewcommand{\arraystretch}{1.52}
\begin{tabular}{|C{4.40cm}|C{4.40cm}|C{4.40cm}|C{4.40cm}|}
  \hline
  \rowcolor{panelgray}
  $y = 6x \Rightarrow a = \blank[1.3cm]$ & $x=4, y=20 \Rightarrow a = \blank[1.1cm]$ & $y=5x,\; x=9 \Rightarrow y = \blank[1.0cm]$ & $y=5x,\; y=45 \Rightarrow x = \blank[1.0cm]$ \\
  \hline
  $\dfrac{y}{x} = 0.8 \Rightarrow y = \blank[1.2cm]$ & $x=10, y=3 \Rightarrow a = \blank[1.1cm]$ & $y=12x,\; x=0 \Rightarrow y = \blank[1.0cm]$ & $y=x+7$៖ សមាមាត្រ? \blank[0.9cm] \\
  \hline
  $x=8, y=-24 \Rightarrow a = \blank[1.0cm]$ & $y=-3x,\; x=6 \Rightarrow y = \blank[0.9cm]$ & $\dfrac{y}{x}=\dfrac{3}{4} \Rightarrow y = \blank[1.2cm]$ & $y=4x,\; y=100 \Rightarrow x = \blank[0.9cm]$ \\
  \hline
  $x=12, y=36 \Rightarrow a = \blank[1.0cm]$ & $y=0.5x,\; x=14 \Rightarrow y = \blank[0.9cm]$ & $y=7x,\; y=63 \Rightarrow x = \blank[0.9cm]$ & $y=2x-1$៖ សមាមាត្រ? \blank[0.8cm] \\
  \hline
\end{tabular}}

\newpage

% ============================================================================
% ទំព័រទី ២៖ E4 លំហាត់គំរូ + E5 ប្រតិបត្តិ + វិនិច្ឆ័យ A-D + Exit Check + កិច្ចការផ្ទះ
% ============================================================================
\secbar{ផ្នែកទី៣៖ ពង្រីកចំណេះដឹងតាមជំហានស៊ើបអង្កេត \enbf{(E4 --} លំហាត់គំរូ ទំ\enbf{.}២៧\enbf{--}២៨\enbf{)}}{រយៈពេល ១០ នាទី}

\textbf{១. លំហាត់គំរូ (ទំ.២៧–២៨)៖} បណ្ណាល័យសាលាផាកពិន័យ $100\text{៛}$ ក្នុង $1$ ថ្ងៃ ចំពោះសៀវភៅសងយឺត ។ តាង $x$ ជាចំនួនថ្ងៃយឺត និង $y$ ជាប្រាក់ពិន័យសរុប (គិតជារៀល)៖

{\renewcommand{\arraystretch}{1.46}
\begin{tabular}{|L{4.50cm}|L{9.10cm}|C{4.20cm}|}
  \hline
  \rowcolor{headgray}
  \textbf{សំណួរ} & \textbf{ការវិភាគតាមមេគុណសមាមាត្រ $a$} & \textbf{ចម្លើយ និងហេតុផល} \\
  \hline
  \textbf{ក.} សរសេរ $y$ ជាអនុគមន៍នៃ $x$ & ក្នុង $1$ ថ្ងៃ ពិន័យ $100\text{៛} \implies$ មេគុណសមាមាត្រ $a = \dfrac{y}{x} = \blank[1.0cm]$ & $y = \blank[1.8cm]$ (\cnum{1}, \cnum{2}) \\
  \hline
  \textbf{ខ.} បើយឺតចំនួន $x = 7$ ថ្ងៃ & ជំនួស $x = 7$ ក្នុងសមីការ៖ $y = 100 \times \blank[0.9cm]$ & $y = \blank[1.6cm]\text{៛}$ (\cnum{3}) \\
  \hline
  \textbf{គ.} បើត្រូវបង់ប្រាក់ $y = 1200\text{៛}$ & ដោះស្រាយ $100x = 1200 \implies x = \dfrac{1200}{\blank[1.0cm]}$ & $x = \blank[1.4cm]$ ថ្ងៃ (\cnum{1}) \\
  \hline
\end{tabular}}

\vspace{4pt}
\textbf{២. ពិនិត្យតារាងតម្លៃ (លំនាំលំហាត់ \#6 ទំ.៤០)៖} តើ $y$ សមាមាត្រស្របនឹង $x$ ដែរឬទេ? បើមែន ចូរសរសេរសមីការ $y=ax$៖

{\renewcommand{\arraystretch}{1.44}
\begin{tabular}{|L{7.20cm}|L{6.20cm}|L{4.40cm}|}
  \hline
  \textbf{ក.} $(x,y)$៖ $(2, 6),\; (4, 12),\; (7, 21),\; (10, 30)$ & ផលធៀប $\frac{y}{x}$៖ $\frac{6}{2}=\blank[0.6cm],\; \frac{12}{4}=\blank[0.6cm],\; \frac{21}{7}=\blank[0.6cm]$ & សមាមាត្រ៖ $y = \blank[1.4cm]$ \\
  \hline
  \textbf{ខ.} $(x,y)$៖ $(1, 4),\; (2, 7),\; (3, 10),\; (4, 13)$ & ផលធៀប $\frac{y}{x}$៖ $\frac{4}{1}=4,\; \frac{7}{2}=\blank[0.8cm]$ (មិនស្មើគ្នា!) & \chkbox\ មែន \quad \chkbox\ មិនមែន (\cnum{4}) \\
  \hline
\end{tabular}}

\vspace{4pt}
$\bullet$ \textbf{សំណួរឆ្លុះបញ្ចាំង \en{(EIA Looking-Back)}៖} ហេតុអ្វីឈ្នួលម៉ូតូឌុប $y = 2000 + 1500x$ មិនមែនជាសមាមាត្រស្រប? \blank[3.4cm]

% --- E5 ---
\secbar{ផ្នែកទី៤៖ វាយតម្លៃលទ្ធផលសិក្សា \enbf{(E5 --} ប្រតិបត្តិ ទំ\enbf{.}២៨ \enbf{+} វិនិច្ឆ័យ \enbf{A--D +} សំបុត្រចេញ\enbf{)}}{រយៈពេល ១០ នាទី}

\textbf{ក. លំហាត់ប្រតិបត្តិ (សៀវភៅពុម្ព ទំ.២៨ លើ និងពង្រឹង \#7 ទំ.៤១ — ធ្វើជាដៃគូ រួចប្ដូរសន្លឹកផ្ទៀង \en{Pause-Review})៖}

{\renewcommand{\arraystretch}{1.48}
\begin{tabular}{|L{8.90cm}|L{8.90cm}|}
  \hline
  \textbf{១) ប្រតិបត្តិ (ទំ.២៨)៖} យន្តហោះហោះល្បឿនថេរ $800\unit{km}$ ក្នុង $1\unit{h}$ ។\par\vspace{2pt}
  $\bullet$ សមីការចម្ងាយ $d$ តាមរយៈពេល $t$៖ $d = \blank[1.8cm]$\par\vspace{2pt}
  $\bullet$ បើហោះ $t = 8\unit{h} \implies d = 800\times \blank[0.8cm] = \blank[1.5cm]\unit{km}$\par\vspace{2pt}
  $\bullet$ បើចម្ងាយ $d = 2000\unit{km} \implies t = \blank[1.5cm]\unit{h}$
  &
  \textbf{២) ពង្រឹង (\#7 ទំ.៤១)៖} គេដឹងថា $y$ សមាមាត្រស្របនឹង $x$ ហើយពេល $x = 4$ គេបាន $y = 12$ ។\par\vspace{2pt}
  $\bullet$ រកមេគុណសមាមាត្រ៖ $a = \dfrac{y}{x} = \dfrac{12}{4} = \blank[1.1cm]$\par\vspace{2pt}
  $\bullet$ សមីការ៖ $y = \blank[1.5cm]$ ; ពេល $x = 9 \Rightarrow y = \blank[1.3cm]$ \\
  \hline
\end{tabular}}

\vspace{5pt}
\textbf{ខ. គូសវង់ជុំវិញជម្រើសត្រឹមត្រូវតែ ១ គត់ ក្នុងចំណោម \en{A, B, C, D} រួចសរសេរលេខហេតុផល \cnum{1}--\cnum{5} នៃជម្រើសខុស៖}

{\renewcommand{\arraystretch}{1.35}
\begin{tabular}{|C{4.40cm}|C{4.40cm}|C{4.40cm}|C{4.40cm}|}
  \hline
  \cellcolor{panelgray}\textbf{ជម្រើស \en{A}} & \cellcolor{panelgray}\textbf{ជម្រើស \en{B}} & \cellcolor{panelgray}\textbf{ជម្រើស \en{C}} & \cellcolor{panelgray}\textbf{ជម្រើស \en{D}} \\
  \hline
  $d=800t \Rightarrow d(8)=6400\unit{km}$ & អាយុបង-ប្អូន $(5,10)\to(10,15)$ & $y = 3x + 5$ សមាមាត្រស្រប & រ៉ឺស័រ $250\unit{g}$ លូត $40\unit{mm}$ \\[3pt]
  \chkbox\ ត្រូវ \quad \chkbox\ ខុស (\blank[0.7cm]) & \chkbox\ ត្រូវ \quad \chkbox\ ខុស (\blank[0.7cm]) & \chkbox\ ត្រូវ \quad \chkbox\ ខុស (\blank[0.7cm]) & \chkbox\ ត្រូវ \quad \chkbox\ ខុស (\blank[0.7cm]) \\
  \hline
\end{tabular}}

\vspace{5pt}
\begin{tcolorbox}[enhanced, sharp corners, colback=white, colframe=black, boxrule=0.85pt, boxsep=2.5pt, top=3.5pt, bottom=3.5pt, left=5pt, right=5pt]
  {\muol\fontsize{8.8pt}{12.5pt}\selectfont សំបុត្រចេញរហ័ស \enbf{(Exit Check --} ៦០ វិនាទី\enbf{)}៖}\\[2pt]
  បើ $y$ សមាមាត្រស្របនឹង $x$ ហើយពេល $x=5$ គេបាន $y=35$៖
  (១) $a = \blank[1.2cm]$ \;\en{|}\;
  (២) សមីការ $y = \blank[1.5cm]$ \;\en{|}\;
  (៣) ពេល $x=8 \Rightarrow y = \blank[1.2cm]$
\end{tcolorbox}

\vspace{5pt}
\begin{tcolorbox}[enhanced, sharp corners, colback=panelgray, colframe=black, boxrule=0.85pt, boxsep=2.5pt, top=3.5pt, bottom=3.5pt, left=5pt, right=5pt]
  {\muol\fontsize{8.8pt}{12.5pt}\selectfont កិច្ចការផ្ទះ \enbf{(}សៀវភៅពុម្ព ទំព័រ ៤០\enbf{--}៤១ លំហាត់ \enbf{\#6(}ក\enbf{,}ខ\enbf{,}គ\enbf{), \#7} និង \enbf{\#10)}៖}\\[2pt]
  $\bullet$ \textbf{លំហាត់ \#6 (ក, ខ, គ ទំ.៤០)៖} ពិនិត្យតារាងតើ $y$ សមាមាត្រស្របនឹង $x$ ឬទេ និងសរសេរសមីការទំនាក់ទំនង ។\\
  $\bullet$ \textbf{លំហាត់ \#7 និង \#10 (ទំ.៤១)៖} គណនាតម្លៃក្នុងទំហំសមាមាត្រស្រប និងអានមុន \textbf{១.២. លក្ខណៈនៃសមាមាត្រស្រប (ទំ.២៨)} ។
\end{tcolorbox}

\end{document}
